\bookchapter{第六十二章：画面づくり——SwiftUI入門}

\chapterintro{画面とは、ユーザーの目に触れるすべてです。SwiftUIはAppleが2019年に発表した宣言型UIフレームワーク——画面を「宣言」するだけで、描画、ダークモード対応、さまざまな画面への対応はシステムが自動で引き受けます。2026年のいま、SwiftUIは新規プロジェクトの圧倒的な主流です。この章では「お絵かき」を学びましょう。}

\booksection{SwiftUIの考え方：宣言型であって命令型ではない}

昔ながらの画面づくりは「命令型」です：システムに「ボタンを作る (20,30)に置く 色を設定する」と、一歩ずつ指示します。SwiftUIは「宣言型」：「画面には『押して』と書かれた赤いボタンが1つある」と直接記述すれば、どう描くか、いつ描き直すかはシステムが面倒を見てくれます。

\begin{lstlisting}[language=Swift]
import SwiftUI

struct ContentView: View {
    var body: some View {      // body：「画面の見た目」を記述する
        VStack {               // 縦に並べる
            Text("こんにちは、世界！")
                .font(.largeTitle)      // 修飾子：大きな見出し
            Button("押して") {
                print("押された")
            }
            .buttonStyle(.borderedProminent)  // 青い塗りつぶしボタン
        }
        .padding()
    }
}
\end{lstlisting}

要点：ビュー＝Viewプロトコルに従う構造体。bodyプロパティが画面を記述し、.モディファイア()でつぎつぎ「メイク」していく。これがSwiftUIの基本のすべてです。

\booksection{基本のビュー：テキスト、ボタン、入力欄、画像}

\begin{center}
\begin{tabularx}{\linewidth}{|X|X|X|}
\hline
\textbf{ビュー} & \textbf{役割} & \textbf{例} \tabularnewline
\hline
Text & 文字を表示 & Text("こんにちは") \tabularnewline
Button & ボタン & Button("押して") \{ 動作 \} \tabularnewline
TextField & 入力欄 & TextField("ヒント", text: \$テキスト) \tabularnewline
Image & 画像を表示 & \texttt{Image(systemName: "star.\allowbreak fill")} \tabularnewline
\hline
\end{tabularx}
\end{center}

\begin{center}\small\color{gray}SwiftUIの4大基本ビュー\end{center}

Image(systemName:)に注目——AppleはSF Symbolsというアイコンライブラリを内蔵していて（数千個の無料アイコン）、“star.fill”のような名前を書くだけで、塗りつぶしの星がそのまま表示できます。これぞiOS開発の隠れた特典：画面づくりでは、ほぼアイコン素材探しが不要です。

\booksection{レイアウトの3点セット：VStack、HStack、ZStack}

\begin{itemize}
  \item VStack：縦に並べる（食堂の行列のように）。
  \item HStack：横に並べる（並んで座るように）。
  \item ZStack：重ねて置く（トランプを積むように。バッジやオーバーレイにぴったり）。
\end{itemize}

\begin{lstlisting}[language=Swift]
VStack(alignment: .leading, spacing: 12) {
    HStack {
        Text("請求書")
        Spacer()                 // 伸びる空間：右端まで押しやる
        Text("-25.5 元").foregroundColor(.red)
    }
    Divider()                    // 区切り線
    Text("メモ：昼ごはん")
}
.padding(16)                     // 四方に余白
.background(Color(.systemGray6)) // 薄いグレーの地
.cornerRadius(12)                // 角丸

// 修飾子の順番に意味がある：先にpadding、後にbackground。背景は余白を含まない
\end{lstlisting}

\begin{tcbwarn}{注意：モディファイアの順番には意味がある}
padding(16).background(グレー)とbackground(グレー).padding(16)は結果が違います——前者は「余白を取ってから色を塗る」（背景に余白は含まれない）、後者は「色を塗ってから余白を取る」（背景に余白が含まれる）ためです。初心者にいちばんよくある「変な白いすき間」は、みんなここから来ています。迷ったら順番を入れ替えてみましょう。プレビューがすぐ答えを教えてくれます。

\end{tcbwarn}

\booksection{状態とバインディング：画面を「生き」させよう}

静的な画面では面白くありません。SwiftUIのルール：画面はつねに状態を忠実に映し出し、状態が変われば画面は自動で描き直されます。

\begin{lstlisting}[language=Swift]
struct CounterView: View {
    @State private var count = 0   // @State：この値が変わったら再描画

    var body: some View {
        VStack {
            Text("\(count)回押しました")
            Button("押して +1") {
                count += 1          // 状態を変えると画面が自動更新
            }
        }
    }
}
\end{lstlisting}

入力欄には「双方向バインディング」を使います——\$マークが「値と書き換え権」をセットで渡してくれます：

\begin{lstlisting}[language=Swift]
struct InputView: View {
    @State private var text = ""

    var body: some View {
        TextField("何か入力", text: $text)  // $による双方向バインディング
            .textFieldStyle(.roundedBorder)
        Text("入力した内容：\(text)")   // 入力中にリアルタイム更新
    }
}
\end{lstlisting}

ユーザー操作

状態が変わる（count += 1）

SwiftUIが変化を検知

画面を自動で再描画（bodyの再計算）

\begin{center}\small\color{gray}図3：SwiftUIの「状態 画面」一方向データフロー\end{center}

\booksection{プレビュー、テーマ、アニメーション：SwiftUIの「裏ワザ」}

\begin{itemize}
  \item ファイルのいちばん下に\#Preview \{ ContentView() \}を足すと、右側のキャンバスがリアルタイムで描画——コードを直すと、見たままがそのまま反映されます。
  \item ダークモードはタダ：色を直書きせず、システムカラー（Color(.systemGray6)、.primary）を使えばダークモードに自動で対応し、1行も追加で書く必要がありません。
  \item アニメーションは1行でOK：withAnimation \{ count += 1 \}とするか、ビューに.animation(.default, value: count)を付ける——切り替えがぬるぬる滑らかに。
\end{itemize}

\begin{lstlisting}[language=Swift]
struct AnimatedHeartView: View {
    @State private var liked = false

    var body: some View {
        Button {
            withAnimation(.spring(duration: 0.4)) {
                liked.toggle()
            }
        } label: {
            Image(systemName: liked ? "heart.fill" : "heart")
                .font(.system(size: 48))
                .foregroundColor(liked ? .red : .gray)
                .scaleEffect(liked ? 1.3 : 1.0)
        }
    }
}
\end{lstlisting}

\begin{tcbinfo}{ワンポイント：まず機能、アニメーションはそのあと}
初心者にありがちな失敗は「画面がまだ正しくできていないのに、アニメーションで悩み始める」こと。まず機能を通るようにして、動きは最後にまとめて足しましょう——アニメーションは華やかさの上乗せ、機能は困ったときの命綱です。これも私たちの一貫した「シンプル第一」の原則です。

\end{tcbinfo}

\begin{tcbpractice}{やってみよう：「気分ボタン」を作る}
\booksection{よく使うモディファイア早見表：この1ページがあれば、画面づくりはもう安心}

\begin{center}
\begin{tabularx}{\linewidth}{|X|X|}
\hline
\textbf{やりたいこと} & \textbf{モディファイア} \tabularnewline
\hline
余白 & \texttt{.\allowbreak padding(16)} \tabularnewline
角丸 & \texttt{.\allowbreak cornerRadius(12)} \tabularnewline
背景色 & \texttt{.\allowbreak background(Color(.\allowbreak systemGray6))} \tabularnewline
フォントと文字サイズ & \texttt{.\allowbreak font(.\allowbreak title.\allowbreak bold())}/\texttt{.\allowbreak font(.\allowbreak system(size: 20))} \tabularnewline
前景色（文字／アイコン） & \texttt{.\allowbreak foregroundColor(.\allowbreak red)} \tabularnewline
幅いっぱいに広げる & \texttt{.\allowbreak frame(maxWidth: .\allowbreak infinity)} \tabularnewline
サイズを固定 & \texttt{.\allowbreak frame(width: 100, height: 44)} \tabularnewline
そろえ方 & \texttt{VStack(alignment: .\allowbreak leading)} \tabularnewline
不透明度 & \texttt{.\allowbreak opacity(0.\allowbreak 6)} \tabularnewline
影 & \texttt{.\allowbreak shadow(radius: 5)} \tabularnewline
タップ操作 & \texttt{.\allowbreak onTapGesture \{ .\allowbreak .\allowbreak .\allowbreak  \}} \tabularnewline
キーボードの種類 & \texttt{.\allowbreak keyboardType(.\allowbreak decimalPad)} \tabularnewline
\hline
\end{tabularx}
\end{center}

\begin{center}\small\color{gray}SwiftUIのよく使うモディファイア早見表\end{center}

覚えられなくても大丈夫——このページは調べるためのもの。実際にいちばん使うのは、上の4つ：padding、角丸、背景、フォントです。

テンプレートプロジェクトを「気分チェンジャー」に改造しましょう：ボタンを1つ用意して、タップするたびにうれしい かなしい おこりの間をぐるぐる切り替え、文字と背景色も連動して変わるようにします。ヒント：3組（文字＋色）を配列に入れ、@Stateにインデックスを保存して、タップのたびに(index + 1) \% 3。できたら友だちに送って遊んでもらい、「え、アプリ作れるの？」という最初のリアクションをもらいましょう。

\end{tcbpractice}

\begin{tcbsummary}{本章のまとめ}
\begin{itemize}
  \item SwiftUIは宣言型：どんな見た目かを記述し、描くのも描き直すのもシステムまかせ
  \item ビュー＝Viewプロトコルに従う構造体。bodyが画面を記述。モディファイアでつぎつぎ「メイク」
  \item 4大基本ビュー：Text / Button / TextField / Image。SF Symbolsのアイコンライブラリは無料でたっぷり
  \item レイアウト3点セット：VStackは縦並び、HStackは横並び、ZStackは重ね置き。Spacerは伸び縮みするすき間
  \item @Stateで画面が生き返る。\$の双方向バインディングで入力欄に値を渡す
  \item \#Previewで見たままを確認。システムカラーならダークモードに自動対応。withAnimationで1行アニメーション
\end{itemize}

\end{tcbsummary}
